\section{Radial selection for integrable densities}
\label{sec:radial-integrable-selection}

The scalar and measure limits in \cref{cor:individual-positive-threshold-exponential} do not by themselves determine a limit potential. For the sublogarithmic densities of \cref{lem:projective-sublogarithmic-density}, the radial flux supplies the missing information. It also distinguishes convergence after subtracting the supremum from convergence after subtracting $\lambda^{-1}\log b$, where $b$ is the subsolution threshold of the density: the latter depends on a first moment of the limiting measure.

Use the product $X=X_0\times Y$, $Y=\mathbb P^k$, $H=h_0+\kappa$, and the probability volume $\mu=V_0^{-1}h_0^m\wedge\kappa^k$ of \cref{sec:zero-lower-mass-product-equations}, with $\dim X_0=m$, $v_-(h_0)=0$ and $k\geq2$. Put $b=\binom NmV_0$, where $N=m+k$. For $0<\alpha<1$, let $U_{\lambda,\alpha}$ be the greatest solution of
\[
 \mathrm{MA}^{\mathrm{ext}}_H(U_{\lambda,\alpha})
       =\mathrm e^{\lambda U_{\lambda,\alpha}}q_\alpha\mu,
 \qquad \lambda>0.
\]
The threshold of this forcing measure is $b>0$, by \cref{thm:zero-lower-mass-product-equation}, so the solutions exist by \cref{cor:individual-positive-threshold-exponential}.

\subsection{Separation and the radial flux}

\begin{lemma}\label{lem:radial-greatest-separated}
There is a $\kappa$-plurisubharmonic function $w_{\lambda,\alpha}$ on $Y$ such that
\[
 U_{\lambda,\alpha}=\lambda^{-1}\log b+w_{\lambda,\alpha},\qquad
 w_{\lambda,\alpha}\geq u_\alpha,
 \qquad
 \mathrm{MA}^{\mathrm{ext}}_\kappa(w_{\lambda,\alpha})
       =\mathrm e^{\lambda w_{\lambda,\alpha}}q_\alpha\kappa^k.
\]
The last measure has mass one. On the affine chart centered at $o$, set $s=\log|z|^2$ and $\rho(s)=\log(1+\mathrm e^s)$, so that $\kappa=\ddc\rho$ there, and write radial functions as functions of $s$. The function $w=w_{\lambda,\alpha}$ is radial and satisfies
\begin{equation}\label{eq:radial-selection-flux}
 \begin{aligned}
 Q_{\lambda,\alpha}(s)
   &=\int_{-\infty}^s\mathrm e^{\lambda w(\sigma)}q_\alpha(\sigma)\,
                                      \mathrm d(\rho'(\sigma)^k),\\
 w'(s)&=Q_{\lambda,\alpha}(s)^{1/k}-\rho'(s),
 \qquad Q_{\lambda,\alpha}(-\infty)=0,
 \quad Q_{\lambda,\alpha}(+\infty)=1.
 \end{aligned}
\end{equation}
\end{lemma}
\begin{proof}
Choose a smooth nondecreasing concave function $\chi$ on $[0,+\infty)$, equal to the identity on $[0,1]$ and constant on $[2,+\infty)$, and put $q_j=j\chi(q_\alpha/j)$. Concavity gives $q_j\uparrow q_\alpha$ and
\[
 \min(q_\alpha,j)\leq q_j\leq\min(q_\alpha,2j).
\]
The functions $q_j$ are smooth and positive on $Y$: near $o$ they are constant, and elsewhere the density is smooth. Let $w_j$ be the continuous solution of $\mathrm{MA}_\kappa(w_j)=\mathrm e^{\lambda w_j}q_j\kappa^k$ from \cref{thm:nguyen-exp}. It is the greatest subsolution by \cref{thm:individual-exponential-subsolution} and bounded uniqueness. Since $u_\alpha\leq0$ and $\mathrm{MA}^{\mathrm{ext}}_\kappa(u_\alpha)=q_\alpha\kappa^k\geq q_j\kappa^k$, the function $u_\alpha$ is a subsolution for every $q_j$, whence $u_\alpha\leq w_j$. Since the $q_j$ increase, \cref{lem:exp-order} shows that the $w_j$ decrease; let $w$ be their limit.

By the product expansion in \cref{thm:zero-lower-mass-product-equation}, the pullback $W_j=\lambda^{-1}\log b+w_j$ is a bounded solution on $X$ of the exponential equation with the smooth density $q_j\mu$ in place of $q_\alpha\mu$. It is therefore also the greatest subsolution of that equation. Every subsolution for $q_\alpha\mu$ is a subsolution for $q_j\mu$, so $U_{\lambda,\alpha}\leq W_j$.

Monotone convergence in capacity and \cref{thm:exponential-closedness} on $Y$ identify the equation for $w\geq u_\alpha$. Moreover, since $\kappa$ is closed, Stokes' theorem gives $\int_Y\mathrm e^{\lambda w_j}q_j\kappa^k=1$. The integrands are dominated by $\mathrm e^{\lambda\sup_Yw_1}q_\alpha$, so dominated convergence gives the same identity for $w$. By the product expansion, the pullback of $\lambda^{-1}\log b+w$ solves the equation on $X$, and hence lies below $U_{\lambda,\alpha}$. Together with $U_{\lambda,\alpha}\leq\lim_jW_j$, this proves separation and all measure assertions.

Uniqueness makes each $w_j$ invariant under the unitary action on the affine coordinates $z\in\mathbb C^k$, and the limit has the same invariance. A radial $\kappa$-plurisubharmonic function has a convex nondecreasing profile $\rho+w$. Symmetric convolution of this convex profile on compact $s$ intervals, decreasing as the scale decreases, and bounded-product continuity give its radial measure $\mathrm d((\rho+w)'^k)$, as in \cref{lem:projective-sublogarithmic-density}. The bound $w\geq u_\alpha$, whose growth is $o(|s|)$ at $-\infty$, forces the initial derivative to vanish. The positive density makes the flux continuous and positive at finite $s$, yielding \eqref{eq:radial-selection-flux}. Its final value is the mass one already proved. The profile is consequently smooth on finite $s$ intervals.
\end{proof}

\subsection{The normalized limit}

Set
\[
 a_{\lambda,\alpha}=\sup_Yw_{\lambda,\alpha},\qquad
 v_{\lambda,\alpha}=w_{\lambda,\alpha}-a_{\lambda,\alpha},\qquad
 c_{\lambda,\alpha}=\mathrm e^{\lambda\sup_XU_{\lambda,\alpha}}
                       =b\mathrm e^{\lambda a_{\lambda,\alpha}}.
\]
In particular $a_{\lambda,\alpha}\geq0$.

\Needspace{12\baselineskip}
\begin{theorem}\label{thm:radial-integrable-selection}\index{Monge--Ampere equation@Monge--Amp\`ere equation!radial selection}\index{Convergence!integrable radial densities}
Let $\alpha_j\to\alpha\in(0,1)$ and $\lambda_j\to0$ with $\lambda_j>0$, at arbitrary relative rates. Then
\[
 v_{\lambda_j,\alpha_j}\longrightarrow u_\alpha
 \quad\text{in }L^1(Y),\text{ in capacity, and locally uniformly on }Y\setminus\{o\}.
\]
Their pullbacks converge in $L^1(X)$, in capacity on $X$, and locally uniformly off $X_0\times\{o\}$. Furthermore $q_{\alpha_j}\to q_\alpha$ in $L^1(\kappa^k)$, $c_{\lambda_j,\alpha_j}\to b$, and the equation measures converge in total variation to $bq_\alpha\mu$. The limiting pullback has finite mass $b$ and belongs to $\mathcal E_{\mathrm{env}}(X,H,0)$.
\end{theorem}
\begin{proof}
Choose $[\alpha_-,\alpha_+]\Subset(0,1)$ containing all late parameters. The density formula in \cref{lem:projective-sublogarithmic-density} gives, uniformly for $\delta$ in this interval and $t=-s$ large,
\[
 q_\delta(s)\leq C\mathrm e^{kt}(1+t)^{-1-k(1-\alpha_+)}.
\]
Away from $o$ the densities are jointly smooth, including on the hyperplane at infinity. Thus one integrable density $g$ dominates all $q_\delta$, and dominated convergence gives $q_{\alpha_j}\to q_\alpha$ in $L^1$.

For every normalized $\kappa$-plurisubharmonic $v$, \cref{thm:normalized-compactness} gives $\int_Y(-v)\kappa^k\leq C_k$. Absolute continuity of the integral of $g$ therefore implies
\[
 \sup_{\sup_Yv=0}\int_{\{v<-T\}}g\kappa^k\longrightarrow0
 \quad(T\to\infty).
\]
Taking $T=\lambda^{-1/2}$ shows that
\[
 \int_Y(1-\mathrm e^{\lambda v})q_\delta\kappa^k
 \leq\lambda T+\int_{\{v<-T\}}g\kappa^k\longrightarrow0
\]
uniformly in $\delta$ and $v$. Write $v_j=v_{\lambda_j,\alpha_j}$ and $a_j=a_{\lambda_j,\alpha_j}$. The exact mass identity of the lemma is
\begin{equation}\label{eq:radial-selection-mass}
 1=\mathrm e^{\lambda_ja_j}
          \int_Y\mathrm e^{\lambda_jv_j}q_{\alpha_j}\kappa^k.
\end{equation}
It gives $\mathrm e^{\lambda_ja_j}\to1$, hence $\lambda_ja_j\to0$. The equation densities $\mathrm e^{\lambda_ja_j}\mathrm e^{\lambda_jv_j}q_{\alpha_j}$ converge in $L^1$ to $q_\alpha$ and are bounded by $2g$ for all late $j$.

Their cumulative radial fluxes $Q_j$ consequently converge uniformly to the flux $Q_\alpha$ of $u_\alpha$. By \eqref{eq:radial-selection-flux}, $v_j'\to u_\alpha'$ uniformly on compact $s$ intervals. The normalized $L^1$ bound controls one value of each profile on a fixed annulus; thus every sequence has a subsequence converging uniformly on compact $s$ intervals to $u_\alpha+d$ for some finite constant $d$.

Near infinity the common bound for the densities is smooth and bounded. Hence $1-Q_j(s)\leq C\mathrm e^{-s}$ and $|v_j'(s)|\leq C\mathrm e^{-s}$ for large $s$. Integrating gives uniform convergence up to the hyperplane at infinity. At the other end, monotonicity of $\rho+v_j$ gives, for $s<s_0$,
\[
 v_j(s)\leq v_j(s_0)+\rho(s_0).
\]
Choose $s_0$ so negative that the right side is negative for all late members of the subsequence. The supremum zero is then attained outside this inner ball. Uniform convergence on its compact complement gives $\sup_Y(u_\alpha+d)=0$, so $d=0$. Normalized compactness identifies every $L^1$ cluster limit with $u_\alpha$; this proves convergence of the entire family and its local uniform convergence off $o$.

For capacity convergence, fix $T>1$. The sublevel $\{u_\alpha<-T\}$ is a neighborhood of $o$, and its complement is a compact set on which convergence is uniform. By \eqref{eq:global-capacity-tail}, the capacity of this sublevel is at most $C/T$. Let $j\to\infty$ and then $T\to\infty$. The same argument on $X$ uses the normalized pullback of $u_\alpha$ and proves convergence in Hermitian capacity there; comparison sets omit pole sets. Pullback also preserves the $L^1$ and local uniform limits. The product expansion proves total-variation convergence of the measures, and \cref{thm:zero-lower-mass-product-equation} supplies the exact limiting mass and envelope class.
\end{proof}

\subsection{Scalar centering and its moment boundary}

\begin{proposition}\label{prop:radial-scalar-centering}\index{Normalization scalar!first moment and centering}
The first moment
\[
 M_\alpha=\int_Yu_\alpha q_\alpha\kappa^k
\]
is finite exactly when $\alpha<k/(k+1)$. In this range,
\begin{equation}\label{eq:radial-centering-constant}
 -M_\alpha=\alpha\int_1^\infty x^{\alpha-1}
 \left[\alpha x^{\alpha-1}+(1-\alpha x^{\alpha-1})\mathrm e^{1-x}\right]^k\,\mathrm dx.
\end{equation}
For the joint limits of \cref{thm:radial-integrable-selection}, if $\alpha<k/(k+1)$, then
\[
 a_{\lambda_j,\alpha_j}\to-M_\alpha,\qquad
 U_{\lambda_j,\alpha_j}-\lambda_j^{-1}\log b
       \to u_\alpha-M_\alpha
\]
in $L^1$, in capacity, and locally uniformly away from the pole fiber. If $\alpha\geq k/(k+1)$, then $a_{\lambda_j,\alpha_j}\to+\infty$ and the latter translated potentials tend to $+\infty$ locally uniformly off the pole fiber, although $\lambda_ja_{\lambda_j,\alpha_j}\to0$. In both cases the exact scalar error is
\[
 \lambda_j^{-1}\log(c_{\lambda_j,\alpha_j}/b)
       =a_{\lambda_j,\alpha_j}.
\]
\end{proposition}
\begin{proof}
As $t=-s\to\infty$, the radial profile of \cref{lem:projective-sublogarithmic-density} satisfies
\[
 u_\alpha\sim-t^\alpha,\qquad
 Q_\alpha'(s)\sim k(1-\alpha)\alpha^k
                         t^{-1-k(1-\alpha)}.
\]
Thus the first moment is finite precisely when $(k+1)\alpha<k$. In that case, write $x=1-\psi$. By \cref{lem:projective-exponential-barrier-growth}, the cumulative flux of $u_\alpha$ is
\[
 Q_\alpha=\left[\alpha x^{\alpha-1}
                  +(1-\alpha x^{\alpha-1})\mathrm e^{1-x}\right]^k.
\]
Integration by parts gives \eqref{eq:radial-centering-constant}; the pole boundary term vanishes since $(k+1)\alpha-k<0$.

Suppose $\alpha<k/(k+1)$ and choose the parameter interval in the preceding proof with $\alpha_+<k/(k+1)$. Its dominating density $g$ satisfies $(-u_{\alpha_+})g\in L^1$. Write $w_j=w_{\lambda_j,\alpha_j}=v_j+a_j$. Since $w_j\geq u_{\alpha_j}\geq u_{\alpha_+}$, the mass identity and $\mathrm e^r\geq1+r$ give $\int_Yw_jq_{\alpha_j}\kappa^k\leq0$. On a fixed annulus $A$ away from $o$, the normalized smooth-volume $L^1$ bound gives
\[
 \int_Aw_jq_{\alpha_j}\kappa^k
       \geq a_j\int_Aq_{\alpha_j}\kappa^k-C.
\]
On the complement, use $w_j\geq u_{\alpha_+}$ and the integrable lower bound. The annulus has uniformly positive measure, so $0\leq a_j\leq C_*$. Along any subsequence with $a_j\to a$, local uniform convergence gives $w_j\to u_\alpha+a$ off $o$. In the exact identity
\[
 \int_Y\frac{\mathrm e^{\lambda_jw_j}-1}{\lambda_j}
                            q_{\alpha_j}\kappa^k=0,
\]
the negative part is dominated by $(-u_{\alpha_+})g$ and the positive part by $C_*\mathrm e^{C_*}g$. Dominated convergence gives $M_\alpha+a=0$. Every subsequence therefore has the asserted scalar limit. Adding these convergent constants to the normalized potentials gives all three claimed topologies.

If $\alpha\geq k/(k+1)$, use instead
\[
 \int_Y\frac{1-\mathrm e^{\lambda_jv_j}}{\lambda_j}
                          q_{\alpha_j}\kappa^k
   =\frac{1-\mathrm e^{-\lambda_ja_j}}{\lambda_j}\leq a_j.
\]
The integrands are nonnegative and tend off $o$ to $(-u_\alpha)q_\alpha$. Fatou's lemma forces $a_j\to+\infty$. The preceding theorem gives $\lambda_ja_j\to0$; local uniform convergence of $v_j$ gives the stated divergence of $w_j$. This includes sequences approaching the moment boundary from below.
\end{proof}

For the elliptic fivefold of \cref{cor:elliptic-fivefold-exact-threshold}, $b=240\pi$ and the moment boundary is $\alpha=2/3$. The choice $\alpha=1/2$ has a finite centering constant. For $2/3\leq\alpha<1$, the supremum-normalized family still converges to the exact threshold solution, while its scalar-centered version diverges off the pole fiber. Thus convergence of the normalization scalars alone does not give a bounded first-order scalar error.

\Needspace{20\baselineskip}
\subsection{Order in the separated class}

We use here the closed-class uniqueness theorem: on a compact connected K\"ahler manifold, two potentials whose non-pluripolar measures have full mass and are equal differ by a constant. This theorem is not proved in this book; its source is given in \cref{hist:radial-integrable-selection}.

\begin{corollary}\label{cor:radial-separated-critical-order}
Every separated threshold subsolution $r=r(y)$ with $r\in\PSH(Y,\kappa)$ and $\mathrm{MA}^{\mathrm{ext}}_H(r)\geq bq_\alpha\mu$ equals $u_\alpha+d$ for some constant $d$. In particular, $u_\alpha$ is the only such subsolution with supremum zero. If $\alpha<k/(k+1)$, then $u_\alpha-M_\alpha$ is the greatest such subsolution whose $q_\alpha$-mean on $Y$ is nonpositive.
\end{corollary}
\begin{proof}
The product formula reduces the inequality to $\mathrm{MA}^{\mathrm{ext}}_\kappa(r)\geq q_\alpha\kappa^k$. Since $\kappa$ is closed, Stokes' theorem gives each bounded truncation mass one, so the extended measure of $r$ has mass at most one. The inequality therefore gives equality and full mass. Closed-class uniqueness yields $r=u_\alpha+d$. In the finite-moment range its mean is $M_\alpha+d$, which is nonpositive exactly when $d\leq-M_\alpha$.
\end{proof}

At and beyond the moment boundary every constant translate has mean $-\infty$, so that mean condition no longer chooses a greatest translate. The order assertion concerns separated potentials. Selection for arbitrary integrable densities and uniqueness of arbitrary critical potentials on the nonclosed product remain the questions of \cref{qu:individual-equation-threshold} and \cref{sec:zero-exponent-topology}.
